Shelfmark 149, scan 29.

\scanpage{29}
\deleted{We may further assume that $\pi_1(X)$ has trivial centre. In characteristic $p>0$, let us note that it suffices to work with the prime-to-$p$ fundamental groups -- and $\unsure{\overline m}$ on the localizations at $\ell\ne p$ of the fundamental groups.}
\mspage{14}

The situation $X,Y,f,g$ comes, by extension of the base field, from an analogous situation over a field $K$ \inserted{an extension} of finite type of $\mathbb Q$. \deleted{Let us choose an} \inserted{Let $\overline K$ be the} algebraic closure \deleted{$\unsure{\overline K|\overline Q}$} of $K$ in $k$, \deleted{\unsure{of the $k$ the}} \inserted{and $\overline X,\overline Y,\overline f,\overline g$ are} obtained from \inserted{$X,Y$,} $f,g$ by extension of the base field from $\unsure{K|Q}$ to $\overline K$. We thus have $\pi_1(\overline f)$, $\pi_1(\overline g)$ satisfying condition d) with in addition $\pi_1(\overline X)$ having trivial centre.

But these hypotheses imply that the extensions $E_{\inserted{X/K}}\deleted{(X)}\inserted{=\pi_1(X)}$, $E_{\inserted{Y/K}}\deleted{(Y)}\inserted{=\pi_1(Y)}$ of $E_{\overline K/K}$ by $\pi_1(\overline X)$, respectively $\pi_1(\overline Y)$, as well as the induced homomorphisms $\pi_1(f)$, $\pi_1(g)$, are reconstructed from \inserted{$\pi_1(\overline X)$, $\pi_1(\overline Y)$, from} $\pi_1(\overline f)$, $\pi_1(\overline g)$ and from the outer action of $E_{\overline K/K}$ on these groups. We shall now show the
\margin{28}
